\documentclass[12pt,a4paper]{article}
\usepackage[T2A]{fontenc}
\usepackage[utf8]{inputenc}
\usepackage[russian]{babel}
\usepackage{geometry}
\geometry{margin=2.5cm}
\usepackage{amsmath,amssymb,mathtools}
\usepackage{enumitem}
\usepackage{hyperref}
\hypersetup{colorlinks=true,linkcolor=black,urlcolor=blue}
\setlength{\parindent}{0pt}
\setlength{\parskip}{6pt}
\begin{document}
{\Large\bfseries
\begin{flushleft}
Отчет по лабораторной работе №1 (Вариант №21)
\end{flushleft}}
\vspace{0.5em}\hrule\vspace{1.2em}

\section*{Задание}
По имеющейся SRS определить:
\begin{itemize}[leftmargin=1.5em]
  \item завершимость
  \item конечность классов эквивалентности по НФ (для построения эквивалентностей считаем, что правила могут применяться в обе стороны). Если их конечное число, то построить минимальную систему переписывания, им соответствующую.
  \item локальную конфлюэнтность и пополняемость по Кнуту--Бендиксу
\end{itemize}

\section*{Исходная система правил (SRS) \(\mathcal{T}\)}
\[
\begin{cases}
aa \to bb\\
ccb \to dca\\
a \to dcc\\
cd \to bc
\end{cases}
\]
Алфавит $\Sigma=\{a,b,c,d\}$.

\section*{Завершимость}
Для доказателсьтва завершимости я использую метод от противного и поочерёдно проанализирую все правила.

\subsection*{2.1 Правило 1 (\(aa \to bb\))}
Рассматриваю применение правила 1 к фрагменту \(w_1aa w_2 \to w_1bb w_2\).

\textbf{A) Возникновение нового \(aa\) в той же позиции.}\\
Пусть \(w_2=a^{k}w_2'\) (\(k\ge0\)). Новое \(aa\) с тем же левым краем появляется только слева за счёт 2 правила:
\[
w_1'ccb\,a^{k}w_2' \to w_1'dc\,a^{k+1}w_2'.
\]
Источником новых \(aa\) на данном стыке может быть только соседнее слева применение 2 правила, и лишь если суффикс заканчивался на \(ccb\).

\textbf{B) Поведение блока \(a^{k+1}\) на стыке.}\\
\(1\) правило укорачивает блок:
\[
w_1'dc\,a^{k+1}w_2' \to w_1'dc\,a^{k-1}bb\,w_2' \to \dots
\]
\(3\) правило на первой \(a\) разрывает стык:
\[
w_1'dc\,a\,a^{k}w_2' \to w_1'dc\,dcc\,a^{k}w_2'.
\]
В результате, число применений 1 правила конечно: блок либо сокращается конечным числом шагов, либо разрывается.

\textbf{C) Пополнения за счёт многократных \(ccb \to dca\) слева.}\\
Для пополнения стыка нужен тот же левый суффикс \(ccb\) с тем же правым концом. После разрыва 3 правилом справа от \(cc\) нет нужной \(a\). :
\[
w_1'dc\,a\,a^{k}w_2' \to w_1'dc\,dcc\,a^{k}w_2',
\]
После свёртки 1 правилом вернуть прежний \(ccb\) на том же правом конце нельзя.:
\[
w_1'ccb\,a\,w_2 \to w_1'dca\,a\,w_2 \to w_1'dc\,bb\,w_2,
\]
Таким образом, один и тот же стык не пополняется бесконечно.

\textbf{Вывод:}
В бесконечной цепочке 1 правило не может срабатывать бесконечно: на каждой фиксированной позиции число его применений конечно. Следовательно, при построении гипотетической бесконечной цепочки 1 правило можно не учитывать.

\subsection*{2.2 Правило 3 (\(a \to dcc\))}

Рассматриваю применение правила 3 к фрагменту \(w_1a w_2 \to w_1dcc w_2\).

\textbf{A) Появление новой \(a\) в той же позиции.}\\
Новая \(a\) появляется только слева по 2 правилу:
\[
w_1' \, ccb \, w_2 \to w_1' \, dca \, w_2.
\]
Если \(w_2 = a^{k} w_2'\) (\(k\ge 0\)), то
\(
w_1' \, ccb \, a^{k} w_2' \to w_1' \, dc \, a^{k+1} w_2'.
\)\\
Источник новой \(a\) на этом стыке - только применение слева 2 правила при суффиксе \(ccb\) у \(w_1\).

\textbf{B) Поведение фрагмента \(dcc\) на стыке.}\\\
4 правило на границе справа при \(w_2 = d\,w_2'\):
\[
w_1 \, dccd \, w_2' \to w_1 \, dcbc \, w_2'.
\]
2 правило справа при \(w_2 = b\,w_2'\):
\[
w_1 \, dccb \, w_2' \to w_1 \, ddca \, w_2'.
\]
1 правило справа при \(w_2 = aa\,w_2'\):
\[
w_1 \, dccaa \, w_2' \to w_1 \, dccbb \, w_2'.
\]
Таким образом, на фиксированном стыке преобразования ограничены по числу шагов

\textbf{C) Пополнения за счёт многократных \(ccb \to dca\) слева.}\\
Чтобы снова получить \(a\) с тем же левым краем после применения правила 3, нужно каждый раз иметь слева суффикс \(ccb\):
\[
w_1' ccb\, a\, w_2 \to w_1' dca\, a\, w_2 \to w_1' dc\, bb\, w_2,
\]
или при \(w_2=a^{k}w_2'\):
\[
w_1' ccb\, a^{k} w_2' \to w_1' dc\, a^{k+1} w_2'.
\]
После преобразований правый край меняется так, что вернуть тот же \(ccb\) на том же правом конце нельзя без новых \(c\) и \(b\), а их появление здесь ограничено описанными ранее преобразованиями.

\textbf{Вывод:}
В бесконечной цепочке 3 правило не может применяться бесконечно: новая \(a\) здесь возникает только один раз, затем преобразования конечны. Таким образом, правило 3 не порождает бесконечность. 

\subsection*{2.3 Правило 2 (\(ccb \to dca\)) и правило 4 (\(cd\to bd\))}
Доказано, что 1 правило (\(aa \to bb\)) и 3 правило (\(a \to dcc\)) не могут участвовать бесконечно в цепочке переписываний. В таком случае цепочка переписываний может быть бесконечна только за счёт применения 2 правила (\(ccb \to dca\)) и 4 правила (\(cd \to bc\)). Проанализирую завершимость системы, состоящей только из этих двух правил.

Введём функцию на словах:
\[
F(w) = 2|w|_c + |w|_d
\]

Проверка правил:\\
- правило \(ccb \to dca\):  
  \(c\) становится на 1 меньше, \(d\) на 1 больше.  \\
  \(F\) меняется так: \(2\cdot(-1)+(+1)=-1\). Значит, \(F\) уменьшается на 1.\\
- правило \(cd \to bc\):  
  \(d\) становится на 1 меньше, \(c\) не меняется.  \\
  \(F\) меняется так: \(0+(-1)=-1\). Значит, \(F\) уменьшается на 1.

\textbf{Вывод:} при каждом шаге по правилам 2 и 4 функция \(F\) уменьшается на 1 и всегда остаётся неотрицательной. Бесконечно так убывать она не может. Значит, система из правил 2 и 4 завершима.

\subsection*{2.4 Итоговый вывод о завершимости SRS}
На основе приведённых доказательств видно, что ни одно из правил не может участвовать бесконечно в цепочке переписываний. Все правила либо приводят к конечному числу применений, либо ограничены в числе шагов. Цепочка переписываний не может быть бесконечно, а следовательно система переписывания является завершимой

\section*{Конечность классов эквивалентности по НФ}
Нормальная форма - это слово, в котором нет ни одного редекса, т.~е. нет ни одной подстроки, совпадающей с левой частью правила переписывания. В нашей системе редексами являются: \(aa\), \(a\), \(cd\) и \(ccb\).\\
Рассмотрим слова вида \(c^n\) при \(n\ge 0\). В них встречается только буква \(c\), поэтому подстроки \(aa\), с, \(cd\) и \(ccb\) там отсутствуют. Следовательно, каждое \(c^n\) уже находится в нормальной форме. Поскольку таких слов бесконечно много, они попарно различны и не могут быть преобразованы друг в друга, множество нормальных форм бесконечно, а значит бесконечно и число классов эквивалентности по НФ.

\section*{Минимальная система переписывания}
Поскольку число классов эквивалентности по НФ бесконечно, минимальную систему переписывания, им соответствующую, построить нельзя.

\section*{Локальная конфлюэнтность}
Для проверки локальной конфлюэнтности рассмотрим слово \(aa\), возникающее при перекрытии правил \(aa\to bb\) и \(a\to dcc\).
К слову \(aa\) применим переписывание тремя способами:
\begin{enumerate}[label=\arabic*) ,leftmargin=1.5em]
  \item по подстроке \(a\)\(a\): \(\underline{a} \to bb\);
  \item через первую \(a\): \(\underline{a}a \to dcc\,\underline{a} \to dc\,\underline{cd}\,cc \to dcbccc\);
  \item через вторую \(a\): \(a\,\underline{a} \to \underline{a}dcc \to dc\,\underline{cd}\,cc \to dcbccc\).
\end{enumerate}

Ветви 2) и 3) сливаются в одно слово и приходят к одной и той же нормальной форме \(dcbccc\).
Ветвь 1) приводит к нормальной форме \(bb\).
Обе формы \(bb\) и \(dcbccc\) неприводимы (к ним нельзя больше применить правила переписывания), и поскольку \(bb \neq dcbccc\), критическая пара \emph{не} замыкается.
Следовательно, данная система не локально конфлюэнтна.

\section*{Пополняемость по Кнуту–Бендиксу}
Введём порядок \emph{shortlex}: сначала сравнение по длине, а при равной длине — лексикографически \((a<b<c<d)\).
Ориентируемая этим порядком система будет иметь имя \(\mathcal{T}'\) и примет следующий вид:
\[
\begin{cases}
bb \to aa\\
dca \to ccb\\
dcc \to a\\
cd \to bc
\end{cases}
\]

Рассмотрим критические пары, возникающие из первых перекрытий исходных правил.\\
Редексы \(cd\) и \(dcc\) имеют переекрытие, в результате чего слово \(cdcc\) можно переписать 2 разными способами:\\
1) \(\underline{cd}cc \to bccc\)\\
2) \(c\underline{dcc} \to ca\)\\
Слова \(bccc\) и \(ca\) не сводятся к общему терму,поэтому для замыкания пары алгоритм добавляет новое правило \(bccc \to ca\) в установленном ранее порядке.

Редексы \(cd\) и \(dca\) имеют переекрытие, в результате чего слово \(cdca\) можно переписать 2 разными способами:\\
1) \(\underline{cd}ca \to bcca\)\\
2) \(c\underline{dca} \to cccb\)\\
Слова \(cccb\) и \(bcca\) не сводятся к общему терму,поэтому для замыкания пары алгоритм добавляет новое правило \(cccb \to bcca\). в установленном ранее порядке..

Далее рассмотрим слово \(cccbccc\), получающееся путем наложения редексов \(cccb\) и \(bccc\) полученных правил, также можно расписать 2 способами:\\
1) \(\underline{cccb}ccc \to bccaccc\)\\
2) \(ccc\underline{bccc} \to cccca\)\\
Слова \(bccaccc\) и \(cccca\) не сводятся к общему терму, поэтому для замыкания пары алгоритм добавляет новое правило \(bccaccc \to cccca\) в установленном ранее порядке.

При наложении редексов \(cccb\) и \(bccaccc\) получится слово \(cccbccaccc\), которое можно расписать следующим образом:\\
1) \(\underline{cccb}ccaccc \to bccaccaccc\)\\
2) \(ccc\underline{bccaccc} \to ccccccca\)\\
Слова \(bccaccaccc\) и \(ccccccca\) не сводятся к общему терму, поэтому для замыкания пары алгоритм добавляет новое правило \(bccaccaccc \to ccccccca\). в установленном ранее порядке.

Наблюдается паттерн: из слова \(cccb(cca)^nccc\), полученного перекрытием \(cccb\) и \(b(cca)^nccc\), возникает правило
\[
bcc(acc)^{\,n+1}c \;\to\; c^{\,3n+4}a \quad \text{для любого } n\ge 0.
\]
Таким образом, алгоритм Кнута-Бендикса будет порождать бесконечную цепочку правил \(bcc(acc)^{\,n+1}c \to c^{\,3n+4}a\) и никогда не завершится, а значит система \(\mathcal{T}'\) не пополняема по Кнуту–Бендиксу.


\section*{Инварианты}
Для данной SRS \(\mathcal{T}\) были найдены следующие инварианты:

1) \textbf{Чётность длины слова.}  \\
\(I(w)=|w|\bmod 2=\big(|w|_{a}+|w|_{b}+|w|_{c}+|w|_{d}\big)\bmod 2.\)

2) \textbf{Паритетный баланс \{a,c\} против \{b,d\}}\\
\(P(w)=\big(|w|_{a}+|w|_{c}-|w|_{b}-|w|_{d}\big)\bmod 2\)

3) \textbf{Матричный инвариант.}  \\
\(M(w)=A^{\;\big(|w|_{a}+|w|_{c}-|w|_{b}-|w|_{d}\big)\bmod 4}.\)

Задаём \(\Psi:\{a,b,c,d\}\to M_{2}(\mathbb{Z})\) (кольцо квадратных матриц \(2\times2\) с целыми элементами):
\[
\begin{gathered}
A=\begin{pmatrix}0&-1\\[2pt]1&0\end{pmatrix},\qquad
B=-A=\begin{pmatrix}0&1\\[2pt]-1&0\end{pmatrix}.
\\[4pt]
\Psi(a)=\Psi(c)=A,\qquad \Psi(b)=\Psi(d)=B.
\end{gathered}
\]
Для слова \(w=x_1x_2\ldots x_k\) полагаем \(M(w)=\Psi(x_1)\Psi(x_2)\cdots\Psi(x_k)\); при этом \(A^{2}=B^{2}=-I\), \(AB=BA=I\). В результате, Для каждой допустимой замены значения на её левой и правой частях совпадают как матрицы; значит, 
\(M(w)\) сохраняется при любом шаге переписывания.

Поскольку \(B=-A\) и \(AB=I\), а матрица \(A\) имеет порядок \(4\) (\(A^{4}=I\)), инвариант \(M(w)\) должен вычисляться как \(A^{\;\big(|w|_{a}+|w|_{c}-|w|_{b}-|w|_{d}\big)\bmod 4}\).

\end{document}
